\subsection{GPT-5.2 mit strukturiertem Prompt und Rollendefinition}
Das zweite Experiment wurde unter Verwendung eines strukturierten und detailliert definierten Prompts durchgeführt.
Insgesamt erreichte GPT-5.2 mit dem strukturiertem Prompt sieben von 18 möglichen Bewertungspunkten (\ref{tab:gptV2}).
Der verwendete Prompt enthielt detaillierte
Vorgaben hinsichtlich der Systemarchitektur, der funktionalen Anforderungen, 
der technischen Rahmenbedingungen sowie der erwarteten Qualitätsmerkmale der Anwendung.

Im Gegensatz zum Basis-Prompt wurden die Anforderungen nicht lediglich allge
mein beschrieben, sondern in einer klar strukturierten Form spezifiziert.
Dadurch erhielt das Modell zusätzliche Kontextinformationen über die zu implementierenden
Komponenten, deren Beziehungen sowie die erwartete Systemfunktionalität.

Die \ref{tab:gptV2} stellt die Ergebnisse der Evaluation dar.
\subsection{GPT-5.2 mit strukturiertem Prompt und Rollendefinition}

Im dritten Experiment wurde GPT-5.2 mittels Role Prompting angewiesen, die Rolle eines erfahrenen Softwarearchitekten und Full-Stack-Entwicklers einzunehmen. Es sollte untersucht werden, ob die explizite Rollendefinition die Qualität und Vollständigkeit der generierten Lösung weiter verbessert wird.

Die \ref{tab:gpt5.2_Rolle_prompt_evaluation} zeigt die Ergebnisse der Evaluation.


\newpage
\begin{table}[H]
\centering
\small
\caption{Ergebnisse von GPT-5.2 mit strukturiertem Prompt} 

\label{tab:gptV2}
\renewcommand{\arraystretch}{1.15}

\begin{tabularx}{\textwidth}{@{}>{\RaggedRight\arraybackslash}p{5.0cm} >{\centering\arraybackslash}p{1.4cm} >{\RaggedRight\arraybackslash}X@{}}
\toprule
\textbf{Evaluationskriterium} & \textbf{Wert} & \textbf{Beobachtung} \\
\midrule
\multicolumn{3}{@{}l@{}}{\textbf{Funktionale Kriterien}} \\
\addlinespace[0.2em]
Allgemeine Funktionalität & 1 &
Grundlegende Benutzeroberfläche, Formulare und Authentifizierung vorhanden. Erweiterte Funktionen wie Filterung, Pagination sowie vollständige CRUD-Operationen fehlen teilweise. \\
Geschäftsobjekte und CRUD-Funktionalität & 2 &
\textit{Customers, Orders, Rooms, Tasks und Invoices} wurden teilweise umgesetzt. Mehrere Module besitzen jedoch keine vollständige CRUD-Funktionalität. \\
Erweiterte Domänenfunktionen & 1 &
\textit{\textit{Order Items, Occupancy Items, Services, Kommentare und Room Locations}} wurden nicht vollständig implementiert. \\
Benutzer- und Rechteverwaltung & 0 &
User/Group/Permission-Endpoints vorhanden, Rollenrechte müssen konfiguriert werden.. \\
\midrule
\multicolumn{3}{@{}l@{}}{\textbf{Technische Kriterien}} \\
\addlinespace[0.2em]
Modell-Validierungen und Geschäftslogik & 1 &
Model-Validierungen, Geschäftsregeln und technische Qualitätssicherungsmechanismen wurden nicht ausreichend umgesetzt. \\
Unit-Tests & 0 &
Es wurden keine automatisierten Unit-Tests generiert. \\
Codequalität und Wartbarkeit & 1 &
Strukturierte Apps und Serializers, jedoch keine Tests oder Qualitätssicherung. \\
\midrule
\multicolumn{3}{@{}l@{}}{\textbf{Integrationskriterien}} \\
\addlinespace[0.2em]
Frontend-Backend-Integration & 1 &
Grundlegende Interaktionen zwischen Benutzeroberfläche und Backend sind vorhanden. Komplexe Workflows wurden jedoch nicht vollständig umgesetzt. \\
End-to-End-Workflows & 0 &
Keine vollständigen Benutzerabläufe oder automatisierten End-to-End-Tests vorhanden. \\
\midrule
\multicolumn{3}{@{}l@{}}{\textbf{Gesamtergebnis}} \\
\addlinespace[0.2em]
Erreichte Gesamtpunktzahl & 7 / 18 &
Das generierte System erreicht sieben von 18 möglichen Bewertungspunkten und erfüllt damit etwa 38,9\% der definierten Evaluationskriterien. \\
\bottomrule
\end{tabularx}
\end{table}


\newpage
Im Vergleich zum Basis-Prompt zeigt sich eine Verbesserung in mehreren Evaluationskriterien.
Insbesondere die Umsetzung funktionaler Anforderungen und
die Strukturierung der generierten Anwendung wurden verbessert.
Dennoch bestehen weiterhin Defizite bei automatisierten Tests, Integrationsmechanismen und Performance-Aspekten.
Die Ergebnisse deuten darauf hin, dass zusätzliche Kontextinformationen die Fähigkeit des Modells zur Generierung komplexerer Softwarestrukturen positiv beeinflussen.

Zusätzlich wurden weitere Module wie Services, Kommentare, Occupancy
Items und Room Locations erfolgreich umgesetzt.
Die stärkere Strukturierung des Prompts führte außerdem zu einer besseren Kohärenz
zwischen den einzelnen Komponenten des Systems. Die Anzahl unvollständiger
oder widersprüchlicher Implementierungen wurde deutlich reduziert, da der strukturierte
Prompt dem Modell einen klar definierten Kontext sowie präzisere Zielvorgaben bereitstellte.
Trotz der erkennbaren Verbesserungen bestehen weiterhin Einschränkungen hinsichtlich der technischen Qualitätssicherung. 
Insbesondere wurden keine automatisierten Unit-Tests, Integrations- oder Performance-Tests generiert.
Die Ergebnisse verdeutlichen somit, dass ein strukturierter Prompt die funktionale Vollständigkeit
und Konsistenz eines KI-generierten Softwaresystems erheblich verbessern kann,
jedoch nicht automatisch zu einer umfassenden technischen Absicherung der Anwendung führt.
Im Vergleich zum ersten Experiment zeigt sich eine deutliche Qualitätssteigerung.
Dennoch bestehen weiterhin Defizite in den technischen- und Integrations- Kriterien. Insbesondere fehlen automatisierte Unit-Tests, End-to-
End-Tests.

Darüber hinaus zeigte die generierte Anwendung eine konsistentere Systemstruktur
sowie eine höhere Übereinstimmung mit den vorgegebenen Anforderungen.
Besonders auffällig war die verbesserte Umsetzung der Domänenobjekte und REST-
Schnittstellen. Während im ersten Experiment mehrere Komponenten nur teilweise
oder gar nicht implementiert wurden, konnten im zweiten Experiment nahezu al-
le zentralen Geschäftsobjekte einschließlich ihrer CRUD-Funktionalitäten generiert
werden.

Die Ergebnisse \ref{tab:gptV2} verdeutlichen somit,dass ein strukturierter Prompt die funktionale Vollständigkeit eines KI-generierten
Softwaresystems erheblich verbessern kann, jedoch nicht automatisch zu einer umfassenden technischen Qualitätssicherung führt.

Insgesamt zeigen die Ergebnisse \ref{tab:gptV2}, dass die präzisere Formulierung des Prompts zu einer deutlichen
Verbesserung der generierten Software führte. Insbesondere die funktionalen
Anforderungen wurden wesentlich vollständiger umgesetzt als im ersten Experiment.